% Generated draft: requires mathematical and source review.
\documentclass[11pt]{article}
\usepackage[margin=25mm]{geometry}
\usepackage{fontspec}
\setmainfont{Latin Modern Roman}
\usepackage{amsmath,amssymb,mathtools,unicode-math}
\setmathfont{Latin Modern Math}
\usepackage{enumitem,needspace,xurl,hyperref}
\setlength{\parindent}{0pt}
\setlength{\parskip}{5pt}
\setlength{\emergencystretch}{3em}
\newcommand{\N}{\mathbb{N}}
\newcommand{\Z}{\mathbb{Z}}
\newcommand{\Q}{\mathbb{Q}}
\newcommand{\R}{\mathbb{R}}
\newcommand{\C}{\mathbb{C}}
\newcommand{\F}{\mathbb{F}}
\newcommand{\A}{\mathbb{A}}
\newcommand{\PP}{\mathbb{P}}
\newcommand{\E}{\mathbb{E}}
\newcommand{\eps}{\varepsilon}
\newcommand{\dd}{\,\mathrm{d}}
\begin{document}
\section*{UnsolvedMath ID: 8700012}
\textbf{Source problem code:} catalogue entry
\subsection{Definitions and mathematical statement}
For a prime \(p\), normalize \(|q|_p=p^{-v_p(q)}\) for nonzero rational \(q\), and \(|0|_p=0\). Use the standard logarithmic rational height \(h(a/b)=\log\max\{|a|,b\}\) for coprime integers \(a,b\), \(b>0\). Logarithms are natural.

The source's Philippon conjecture asks for global constants
\[
 0<\varepsilon<\tfrac12,\qquad\alpha\ge1,\qquad\beta\ge0,
 \qquad B\in\mathbb Z_{>0},
\]
such that for every \(x,y\in\mathbb Q^\times\) with \(xy^B\ne1\), on putting
\[
 q=xy^B+1,\qquad S=\{p\text{ prime}:|q|_p<1\},
\]
one has
\[
 -\sum_{p\in S}\log|xy^B+1|_p
 \le B\left(\alpha h(x)+\varepsilon h(y)
 +(\alpha B+\varepsilon)\left(\beta+\sum_{p\in S}\log p\right)\right).
\]
The same \(\varepsilon,\alpha,\beta,B\) must work for all admitted rational pairs. When \(q\ne0\), \(S\) is finite: it consists of primes dividing the numerator of \(q\) in lowest terms.

There is a material sign/domain problem in the supplied statement. Its condition \(xy^B\ne1\) does not exclude \(xy^B=-1\); then \(q=0\), every prime lies in \(S\), and the displayed finite estimate involves infinite sums and \(\log0\). One possible correction would exclude \(xy^B=-1\); another would change the expression to \(xy^B-1\). They are different formulations, and neither is adopted without the original source. Thus the inequality is retained exactly as printed with this unresolved exceptional case, not certified as a well-posed universal conjecture.

The research summary describes an open quantitative \(p\)-adic/S-unit estimate related to gcd, \(abc\), and logarithmic-form problems, but supplies no exact proof or auxiliary references. The source's 2004 open presentation and later unresolved assessment are dated attributions pending a primary-literature audit.
\subsection{Short English statement}
Seek universal constants bounding a sum of \(p\)-adic valuations of \(xy^B+1\) by heights and the primes involved. The printed exclusion has a sign mismatch: it permits \(xy^B=-1\), where the displayed sums are not finite. That source issue must be resolved before the conjecture is well posed.
\subsection{Sources}
\begin{itemize}
\item Local registry: \url{data/unsolvedmath/problems.json}, permanent UnsolvedMath ID \(8700012\), source code catalogue entry; supplied statement, background, and research summary, updated 2026-08-06.
\item Michel Waldschmidt, \emph{Open Diophantine Problems} (2004), Conjecture 2.4, attributed to Philippon: \url{https://arxiv.org/abs/math/0312440}. Supplied source-TeX extraction accessed 2026-07-29; a current primary-literature status audit remains required.
\item The summary mentions \(p\)-adic logarithmic estimates, S-units, gcd questions and \(abc\) without separate precise bibliographic references.
\end{itemize}
\end{document}
